Sigui el sistema d'equacions lineals següent, que depèn del paràmetre real $\lambda$:
\[
\left\{\begin{aligned}
x+2\lambda y+(2+\lambda)z&=0\\
(2+\lambda)x+y+2\lambda z&=3\\
2\lambda x+(2+\lambda)y+z&=-3
\end{aligned}\right.
\]

\begin{apartats}

\apartat{1,25}
Discutiu el sistema per als diferents valors del paràmetre $\lambda$.

\begin{solucio}
La matriu del sistema i la matriu ampliada són
$M|\overline{M}=\left(\begin{array}{ccc|c}1&2\lambda&2+\lambda&0\\2+\lambda&1&2\lambda&3\\2\lambda&2+\lambda&1&-3\end{array}\right)$.
Calculem el determinant de la matriu del sistema i estudiem per a quins valors s'anul·la:
\[
\begin{vmatrix}1&2\lambda&2+\lambda\\2+\lambda&1&2\lambda\\2\lambda&2+\lambda&1\end{vmatrix}=9\lambda^3+9=9\left(\lambda^3+1\right).
\]
Quan resolem l'equació $9\left(\lambda^3+1\right)=0$, obtenim $\lambda^3=-1$ i, d'aquí, $\lambda=-1$. Per tant, tenim dos
casos per discutir.\\
\textbf{Cas I: $\lambda\neq-1$.} Com que $|M|\neq0$, aleshores $\operatorname{rang}M=3=\operatorname{rang}\overline{M}=$ nombre
d'incògnites. En virtut del teorema de Rouché-Frobenius, es tracta d'un \fbox{sistema compatible determinat (SCD)} i, per
tant, el sistema té una única solució.\\
\textbf{Cas II: $\lambda=-1$.}
$M|\overline{M}=\left(\begin{array}{ccc|c}1&-2&1&0\\1&1&-2&3\\-2&1&1&-3\end{array}\right)$. Com que $|M|=0$, aleshores
$\operatorname{rang}M<3$, però, com que $\begin{vmatrix}1&-2\\1&1\end{vmatrix}=1+2=3\neq0$, $\operatorname{rang}M=2$. D'altra
banda, si orlem l'anterior menor,
$\begin{vmatrix}1&-2&0\\1&1&3\\-2&1&-3\end{vmatrix}=-3+12-3-6=0$ i, per tant, $\operatorname{rang}\overline{M}=2$. Aleshores,
com que $\operatorname{rang}M=2=\operatorname{rang}\overline{M}<$ nombre d'incògnites $=3$, es tracta d'un \fbox{sistema compatible
indeterminat (SCI)} amb $3-2=1$ grau de llibertat.\\
De forma alternativa, el primer determinant es podria haver calculat utilitzant les propietats dels determinants,
sumant a la tercera fila les dues primeres:
\[
\begin{vmatrix}1&2\lambda&2+\lambda\\2+\lambda&1&2\lambda\\2\lambda&2+\lambda&1\end{vmatrix}
=(3+3\lambda)\begin{vmatrix}1&2\lambda&2+\lambda\\2+\lambda&1&2\lambda\\1&1&1\end{vmatrix}
=3(1+\lambda)\left(3\lambda^2-3\lambda+3\right)=9(1+\lambda)\left(\lambda^2-\lambda+1\right).
\]

\textit{Pauta oficial:} 0,25 per l'expressió de les matrius del sistema, 0,25 pel càlcul del determinant, 0,25 per
tenir el punt singular que anul·la el determinant, 0,25 per la discussió del cas $\lambda\neq-1$ i 0,25 per la
discussió del cas $\lambda=-1$.
\end{solucio}

\apartat{1,25}
Per al cas $\lambda=-1$, resoleu el sistema, interpreteu-lo geomètricament i identifiqueu-ne la solució.

\begin{solucio}
Sabem que es tracta d'un SCI amb un grau de llibertat, i el menor d'ordre 2 no nul indica les equacions i les
incògnites que ens hem de quedar: $\left\{\begin{aligned}x-2y&=-z\\x+y&=3+2z\end{aligned}\right.$. Ara podem resoldre
aplicant el mètode de Cramer:
\[
x=\frac{\begin{vmatrix}-z&-2\\3+2z&1\end{vmatrix}}{\begin{vmatrix}1&-2\\1&1\end{vmatrix}}=\frac{3z+6}{3}=z+2,\qquad
y=\frac{\begin{vmatrix}1&-z\\1&3+2z\end{vmatrix}}{\begin{vmatrix}1&-2\\1&1\end{vmatrix}}=\frac{3+3z}{3}=z+1 .
\]
Solució: \fbox{$x=z+2$, $y=z+1$, amb $z$ com a paràmetre}.\\
Interpretació geomètrica: cadascuna de les equacions és un pla de $\mathbb{R}^3$. El fet que la solució tingui un grau de
llibertat ens indica que aquests tres plans es tallen en una recta, que és la formada pels punts de la forma
$(z+2,\,z+1,\,z)$. Com que $(z+2,\,z+1,\,z)=(2,1,0)+z(1,1,1)$, la recta intersecció és la que passa pel punt $(2,1,0)$
i té vector director $(1,1,1)$.

\textit{Pauta oficial:} 0,25 per la identificació del tipus de solució, 0,25 per la resolució de les incògnites $x$ i
$y$ (o les que escaigui), 0,25 per l'expressió de la solució del sistema, 0,25 per indicar que els tres plans es tallen
en una recta i 0,25 per identificar la recta solució.
\end{solucio}

\end{apartats}
